\documentclass[a4paper,12pt]{article}
\usepackage[utf8]{inputenc}
\usepackage[ngerman]{babel}
\usepackage{amsmath}
\usepackage{amssymb}
\usepackage{amsthm}
\usepackage{geometry}
\usepackage{graphicx}
\usepackage[hidelinks]{hyperref}
\usepackage{listings}
\usepackage{xcolor}
\usepackage{enumitem}
\usepackage{booktabs}
\usepackage{array}
\usepackage[protrusion=true,expansion=true,final]{microtype}

\geometry{margin=2.5cm}
\tolerance=2000
\emergencystretch=3em
\hfuzz=0.5pt
\vfuzz=0.5pt

\renewcommand{\qedsymbol}{$\blacksquare$}

\theoremstyle{plain}
\newtheorem{satz}{Satz}
\newtheorem{lemma}{Lemma}
\newtheorem{korollar}{Korollar}
\theoremstyle{definition}
\newtheorem{definition}{Definition}
\newtheorem{beispiel}{Beispiel}
\theoremstyle{remark}
\newtheorem{bemerkung}{Bemerkung}

\setlist[itemize,1]{label=--}
\definecolor{mainblue}{RGB}{25,70,140}
\definecolor{accentred}{RGB}{180,50,30}
\hypersetup{colorlinks=true,linkcolor=mainblue,citecolor=accentred,urlcolor=mainblue}
\lstset{basicstyle=\ttfamily\small,keywordstyle=\color{blue},
  commentstyle=\color{gray},numbers=left,numberstyle=\tiny,
  breaklines=true,frame=single,
  literate={Ö}{{\"O}}1{Ä}{{\"A}}1{Ü}{{\"U}}1
           {ß}{{\ss}}1{ü}{{\"u}}1{ä}{{\"a}}1{ö}{{\"o}}1}

\newcommand{\ITM}{\mathcal{T}_{\mathrm{ion}}}
\newcommand{\IRF}{\mathcal{F}_{\mathrm{ion}}}
\newcommand{\Nat}{\mathbb{N}}
\newcommand{\Real}{\mathbb{R}}
\newcommand{\Comp}{\mathbb{C}}
\newcommand{\cth}{c_{\mathrm{th}}}
\newcommand{\chigh}{c_{\mathrm{high}}}
\newcommand{\clow}{c_{\mathrm{low}}}
\newcommand{\kth}{\kappa_{\mathrm{th}}}
\newcommand{\kges}{\kappa_{\mathrm{ges}}}

\title{Ionotronische Turing-Maschine:\\[6pt]
Wasserbasierte Berechnungsuniversalit\"at,\\
Ionisches Fl\"ussigkeits-Rechnen (IRF)\\
und photonisch-ionische Synergie (ionoxoqtum)}
\author{Stephan Epp\\[4pt]\small Universit\"at Bielefeld}
\date{8.\ Juni 2026}

\begin{document}
\maketitle

\begin{abstract}
Diese Arbeit entwickelt eine vollst\"andige formale Theorie des
\emph{Ionotronischen Fl\"ussigkeits-Rechnens} (IRF) und beweist,
dass ein wasserbasiertes Medium dasselbe Berechnungsuniversum wie eine
klassische Turing-Maschine realisiert.
Kernresultate:
(i)~\emph{Kohlrausch-Schaltungssatz} (Satz~\ref{satz:kohlrausch}):
Alle Booleschen Grundgatter (AND, OR, NOT) sind ionisch realisierbar;
Beweis durch explizite Schwellenanalyse der Kohlrausch-Leitf\"ahigkeit.
(ii)~\emph{Lemma der ionischen Injektivit\"at} (Lemma~\ref{lem:inject}):
Die Konzentrations-Bit-Kodierung $\mathrm{enc}: \{0,1\} \to \Gamma$ ist
bijektiv und stabil gegen\"uber thermischem Rauschen oberhalb $\cth$.
(iii)~\emph{Physikalische Realisierbarkeit} (Satz~\ref{satz:phys_real}):
Jeder ITM-\"Ubergangsschritt ist mikrofluidisch implementierbar mit
Zeitkomplexit\"at $O(1)$ pro Schritt.
(iv)~\emph{Funktionale Vollst\"andigkeit} (Satz~\ref{satz:vollstaendig}):
Das Gatter-System $\{\mathrm{AND},\mathrm{NOT}\}$ ist funktional vollst\"andig
(de~Morgan-Beweis).
(v)~\emph{Analog-Informationsgehalt} (Satz~\ref{satz:analog_info}):
Ionische Kan\"ale tragen $I = \log_2(\mathrm{DR})\cdot\ell$ bit$\cdot$m
an analoger Information.
(vi)~\emph{Hybrid-Satz} (Satz~\ref{satz:hybrid}):
Plasmon-Resonanzwellenl\"ange $\Delta\lambda = S\cdot(dn/dc)\cdot c$
erlaubt kontaktloses optisches Auslesen.
(vii)~\emph{ITM-Universalit\"atssatz} (Satz~\ref{satz:itm_universal}):
ITM $\equiv_T$ TM (Turing-\"Aquivalenz) durch explizite Simulations-Konstruktion.
Alle Beweise sind vollst\"andig formal und self-contained.
\end{abstract}

\tableofcontents
\newpage

%%%%%%%%%%%%%%%%%%%%%%%%%%%%%%%%%%%
\section{Einf\"uhrung}
%%%%%%%%%%%%%%%%%%%%%%%%%%%%%%%%%%%

\subsection{Motivation und wissenschaftliche Einordnung}

Das Silizium-Transistor-Paradigma, das die Computerarchitektur seit
Shockley, Bardeen und Brattain (1947)~\cite{shockley1948} dominiert,
n\"ahert sich seinen fundamentalen Grenzen.
Die Dennard-Skalierung~\cite{dennard1974} -- die Reduktion von Transistor-Dimensionen
bei konstanter Leistungsdichte -- versagt seit den 2000er-Jahren,
da Leckstr\"ome und Quantentunnelung bei Gate-Breiten unter $5\,$nm
dominieren~\cite{haensch2006}.
Das thermische Plateau liegt bei $\approx 100\,$W/cm$^2$; f\"ur
Exascale-Computing~\cite{dongarra2011} m\"ussen grundlegend andere
Rechenparadigmen erschlossen werden.

Ionotronik -- die Nutzung ionischer Tr\"ager (Na$^+$, K$^+$, Cl$^-$
in w\"assriger L\"osung) statt Elektronen -- bietet eine biologisch
fundierte Alternative~\cite{yang2018,robin2021,bocquet2020}.
Das menschliche Gehirn erzielt mit $\approx 20\,$W eine Rechenleistung,
die aktuelle Supercomputer bez\"uglich Energie-Effizienz um Gr\"o{\ss}enordnungen
\"ubertrifft~\cite{markram2006}.
Die Hodgkin-Huxley-Gleichungen~\cite{hodgkin1952} beschreiben dieses
ionotronische Rechenprinzip exakt: Ionenfl\"usse durch spannungsgesteuerte
Membrankan\"ale kodieren, verarbeiten und \"ubertragen Information.

Diese Arbeit formalisiert Ionotronik als vollst\"andiges \emph{Berechnungsparadigma}
im Sinne der Berechenbarkeitstheorie~\cite{sipser2013}.
Wir konstruieren eine \emph{Ionotronische Turing-Maschine} (ITM) als physikalische
Instantiierung eines abstrakten Rechenmodells, beweisen ihre Turing-\"Aquivalenz
und verbinden sie mit dem Analog-Paradigma~\cite{epp2026analog} sowie
der photonischen Nanooptik~\cite{epp2026oqtum}.

\subsection{Grundlagen der Ionotronik}

\subsubsection{Ionentransport in w\"assrigen L\"osungen}

W\"assrige Elektrolytl\"osungen bestehen aus polaren Wassermolek\"ulen
(Dipolmoment $\mu = 1{,}85\,$D) und dissoziierten Ionen.
F\"ur NaCl:
\[
\mathrm{NaCl} \xrightarrow{H_2O} \mathrm{Na}^+ + \mathrm{Cl}^-.
\]
Die Ionen bewegen sich im elektrischen Feld $\mathbf{E}$ mit Driftgeschwindigkeit
$\mathbf{v}_d = \mu_{\mathrm{ion}} \mathbf{E}$, wobei $\mu_{\mathrm{ion}}$
die ionische Mobilit\"at (m$^2$/(V$\cdot$s)) ist.
Der Ionenstrom: $\mathbf{J} = \sigma \mathbf{E}$, mit Leitf\"ahigkeit $\sigma = \kappa$.

\subsubsection{Abgrenzung zum Elektronenstrom}

Elektronenstrom in Metallen: Tr\"ager sind Elektronen mit Masse
$m_e = 9{,}109 \cdot 10^{-31}\,$kg und Ladung $-e$.
Ionenstrom: Tr\"ager sind Ionen mit Masse $m_{\mathrm{ion}} \approx 10^{-26}$--$10^{-25}\,$kg
und Ladung $\pm ze$.
Konsequenz: Ionentransport ist $\approx 10^5$ mal langsamer als Elektronentransport
($v_{\mathrm{drift}}^{\mathrm{ion}} \approx 10^{-9}$\,m/s bei $E = 100\,$V/m
vs.\ $v_{\mathrm{drift}}^e \approx 10^{-4}$\,m/s).
Daf\"ur ist der Energieverbrauch pro \"ubertragener Ladungseinheit minimal:
$E_{\mathrm{ion}} \approx k_B T \approx 4 \cdot 10^{-21}\,$J (thermisch).

\subsection{Verwandte Arbeiten}

Yang et al.~\cite{yang2018} demonstrierten ionische Gleichrichter und
memristives Verhalten in weichen ionotronischen Kan\"alen aus
PEG-Hydrogel.
Robin et al.~\cite{robin2021} zeigten Langzeitged\"achtnis und
Synapse-artiges Verhalten in 2D-Nanofluidik-Kan\"alen (MoS$_2$).
Bocquet~\cite{bocquet2020} \"ubersieht den Stand der Nanofluidik;
Ionenkan\"ale von $1\,$nm Durchmesser zeigen Quanteneffekte.
Die vorliegende Arbeit ist die erste, die \emph{Turing-Vollst\"andigkeit}
f\"ur ionotronische Systeme formal beweist.

\subsection{Einordnung in das Forschungsportfolio}

Das Ionotronische Fl\"ussigkeits-Rechnen (iono~\cite{epp2026iono}) liefert
das Rechenmedium.
Die aktive Nanooptik (oqtum~\cite{epp2026oqtum}) liefert die photonische
Ausleseschicht.
Das Hybridprinzip ionoxoqtum~\cite{epp2026ionoxoqtum} vereint beide.
Das Analog-Paradigma~\cite{epp2026analog}: Wasser als prim\"ar kontinuierliches
Medium verkn\"upft sich mit der Diskretisierung~\cite{epp2026digi}.
Die neuen Forschungsfelder NF-2 und NF-6~\cite{epp2026nf_all} werden
durch diese Arbeit formal fundiert.
Der Subgraph Algorithmus~\cite{epp2026subgraph} erm\"oglicht polynomielle
Mustererkennung auf ionischen Bandinhalten.

\subsection{Beitr\"age dieser Arbeit im \"Uberblick}

\begin{enumerate}
\item Formale Definition der ITM als 7-Tupel
      mit vollst\"andiger physikalischer Interpretation.
\item Vollst\"andige Beweise aller Logikgatter aus der Kohlrausch-Gleichung.
\item Beweis der Turing-\"Aquivalenz via explizite Simulation.
\item Formalisierung des ionisch-photonischen Hybrids (ionoxoqtum).
\item Verbindung zum Analog-Paradigma und Shannon-Theorie.
\item Ausf\"uhrliche Beweise aller Lemmata und Zwischenergebnisse.
\end{enumerate}

%%%%%%%%%%%%%%%%%%%%%%%%%%%%%%%%%%%
\section{Physikalische Grundlagen des Ionotronischen Rechnens}
%%%%%%%%%%%%%%%%%%%%%%%%%%%%%%%%%%%

\subsection{Thermodynamik der Ionenl\"osungen}

\begin{definition}[Elektrolytl\"osung]
\label{def:elektrolyt}
Eine \emph{Elektrolytl\"osung} der Konzentration $c$ [mol/L] ist eine
wässrige L\"osung vollst\"andig dissoziierten Salzes.
F\"ur ein 1:1-Elektrolyt (z.~B.\ NaCl) gilt:
\[
[\mathrm{Na}^+] = [\mathrm{Cl}^-] = c,
\quad [\mathrm{H}^+][\mathrm{OH}^-] = K_w = 10^{-14}\,(\mathrm{mol/L})^2.
\]
Die \emph{Ionenst\"arke} der L\"osung:
$I = \frac{1}{2}\sum_i z_i^2 c_i = c$ f\"ur 1:1-Elektrolyte.
\end{definition}

\begin{definition}[Ionische Leitf\"ahigkeit nach Kohlrausch]
\label{def:leitfaehigkeit}
Die \emph{molare Leitf\"ahigkeit} $\Lambda_m$ einer Elektrolytl\"osung
der Konzentration $c$ erf\"ullt das \emph{Kohlrausch-Gesetz}:
\[
\Lambda_m(c) = \Lambda_m^\infty - K_{\mathrm{K}}\,\sqrt{c},
\]
wobei $\Lambda_m^\infty$ die \emph{Grenzleitf\"ahigkeit} (bei $c \to 0$,
d.~h.\ unendlich verd\"unnt) und $K_{\mathrm{K}} > 0$ der
\emph{Kohlrausch-Koeffizient} ist.
Die \emph{spezifische Leitf\"ahigkeit} ist dann:
\[
\kappa(c) = \Lambda_m(c) \cdot c = \left(\Lambda_m^\infty - K_{\mathrm{K}}\sqrt{c}\right) c.
\]
F\"ur NaCl bei $T = 25\,^\circ$C:
$\Lambda_m^\infty = 126{,}4 \cdot 10^{-4}\,\mathrm{S\,m^2\,mol^{-1}}$,
$K_{\mathrm{K}} = 60{,}2 \cdot 10^{-4}\,\mathrm{S\,m^2\,mol^{-3/2}\,L^{1/2}}$.
\end{definition}

\begin{lemma}[Monotonie und Unimodalit\"at von $\kappa(c)$]
\label{lem:kappa_mono}
Die Funktion $\kappa: [0, \infty) \to [0, \infty)$, $c \mapsto (\Lambda_m^\infty - K_{\mathrm{K}}\sqrt{c})c$,
hat folgende Eigenschaften:
\begin{enumerate}[label=(\roman*)]
    \item $\kappa(0) = 0$,
    \item $\kappa$ ist auf $[0, c^*]$ streng monoton wachsend und auf $[c^*, \infty)$ streng monoton fallend,
          wobei $c^* = \left(\frac{2\Lambda_m^\infty}{3K_{\mathrm{K}}}\right)^2$,
    \item F\"ur den Betrieb bei $c \leq c^*/4$ ist $\kappa$ streng monoton wachsend.
\end{enumerate}
\end{lemma}

\begin{proof}
Berechne die Ableitung von $\kappa$ bez\"uglich $c$:
\[
\kappa'(c) = \frac{d}{dc}\!\left[(\Lambda_m^\infty - K_{\mathrm{K}}\sqrt{c})\,c\right]
= \Lambda_m^\infty - K_{\mathrm{K}}\sqrt{c} - \frac{K_{\mathrm{K}}}{2\sqrt{c}}\cdot c
= \Lambda_m^\infty - \frac{3}{2}K_{\mathrm{K}}\sqrt{c}.
\]
(i) Einsetzen $c=0$: $\kappa(0) = (\Lambda_m^\infty - 0)\cdot 0 = 0$. \checkmark

(ii) $\kappa'(c) = 0 \Leftrightarrow \Lambda_m^\infty = \frac{3}{2}K_{\mathrm{K}}\sqrt{c^*}
\Leftrightarrow \sqrt{c^*} = \frac{2\Lambda_m^\infty}{3K_{\mathrm{K}}}
\Leftrightarrow c^* = \left(\frac{2\Lambda_m^\infty}{3K_{\mathrm{K}}}\right)^2$.
F\"ur $c < c^*$: $\kappa'(c) > 0$ (wachsend).
F\"ur $c > c^*$: $\kappa'(c) < 0$ (fallend).
Da $\kappa''(c) = -\frac{3K_{\mathrm{K}}}{4\sqrt{c}} < 0$ f\"ur alle $c > 0$,
ist $c^*$ ein globales Maximum.

(iii) F\"ur $c \leq c^*/4 < c^*$ gilt Monotonie nach (ii). \checkmark
\end{proof}

\begin{bemerkung}
F\"ur NaCl: $c^* = (2\cdot 126{,}4/(3\cdot 60{,}2))^2 \approx (0{,}1402)^2 \cdot 10^0
\approx 1{,}96 \cdot 10^{-2}\,\mathrm{mol/L}\cdot (10^0)^2$.
Numerisch: $c^* \approx 1{,}96\,\mathrm{mol/L}$.
Im ionotronischen Betrieb wählen wir $c_{\mathrm{logic}} = 0{,}1\,\mathrm{mol/L}
\ll c^*$, sodass Monotonie gew\"ahrleistet ist.
\end{bemerkung}

\subsection{Das Nernst-Potential}

\begin{definition}[Nernst-Potential]
\label{def:nernst}
Sei $\mathcal{M}$ eine selektiv-permeable Membran, die nur Ionen der
Ladungszahl $z$ passieren l\"asst.
Das \emph{Nernst-Gleichgewichtspotential} bei Innenkonzentration
$c_{\mathrm{in}}$ und Au{\ss}enkonzentration $c_{\mathrm{out}}$ ist:
\[
E_{\mathrm{Nernst}} = \frac{RT}{zF}\ln\frac{c_{\mathrm{out}}}{c_{\mathrm{in}}},
\]
wobei $R = 8{,}3145\,\mathrm{J\,mol^{-1}\,K^{-1}}$, $T$ die absolute Temperatur
und $F = 96\,485\,\mathrm{C\,mol^{-1}}$ die Faraday-Konstante.
\end{definition}

\begin{lemma}[Nernst-Stetigkeit und Monotonie]
\label{lem:nernst_prop}
Die Funktion $E_{\mathrm{Nernst}}: (0,\infty) \to \mathbb{R}$,
$c_{\mathrm{in}} \mapsto E_{\mathrm{Nernst}}(c_{\mathrm{in}})$
(bei festem $c_{\mathrm{out}}$), ist
\begin{enumerate}[label=(\roman*)]
    \item stetig und streng monoton fallend in $c_{\mathrm{in}}$,
    \item surjektiv auf $\mathbb{R}$,
    \item $C^\infty$ f\"ur alle $c_{\mathrm{in}} > 0$.
\end{enumerate}
\end{lemma}

\begin{proof}
(i) $\frac{dE}{dc_{\mathrm{in}}} = \frac{RT}{zF} \cdot \frac{1}{\ln 10} \cdot \frac{-1}{c_{\mathrm{in}}} < 0$
f\"ur $z > 0$ und $c_{\mathrm{in}} > 0$. Stetigkeit folgt aus der Stetigkeit
des Logarithmus auf $(0, \infty)$.

(ii) $\lim_{c_{\mathrm{in}} \to 0^+} E = +\infty$ und
$\lim_{c_{\mathrm{in}} \to \infty} E = -\infty$ (f\"ur $z > 0$).
Zwischenwertsatz liefert Surjektivit\"at.

(iii) $\ln$ ist auf $(0, \infty)$ glatt ($C^\infty$). \checkmark
\end{proof}

\begin{figure}[h!]
\centering
\includegraphics[width=0.95\textwidth]{plot1_iono_schaltung.pdf}
\caption{Links: Ionotronischer Schaltkreis (IRF) in Analogie zur CMOS-Logik.
Farbige Mikrokanalschichten kodieren verschiedene Funktionen.
Rechts: Ionische Leitf\"ahigkeit $\kappa(c)$ nach Kohlrausch f\"ur NaCl
und KCl; Betriebspunkte sind annotiert.}
\label{fig:iono_schaltung}
\end{figure}

\subsection{Elektrochemisches Potential und Donnan-Gleichgewicht}

\begin{definition}[Elektrochemisches Potential]
\label{def:elektrochem}
Das \emph{elektrochemische Potential} der Ionenspezies $i$ in Phase $\alpha$:
\[
\tilde{\mu}_i^\alpha = \mu_i^{0,\alpha} + RT\ln a_i^\alpha + z_i F \phi^\alpha,
\]
wobei $\mu_i^{0,\alpha}$ das Standardpotential, $a_i^\alpha$ die Aktivit\"at
(bei kleinen Konzentrationen $a_i \approx c_i$) und $\phi^\alpha$ das
elektrische Potential in Phase $\alpha$.
\end{definition}

\begin{lemma}[Donnan-Gleichgewicht]
\label{lem:donnan}
Im thermodynamischen Gleichgewicht zwischen zwei Phasen $\alpha$ und $\beta$,
die durch eine semipermeable Membran getrennt sind, gilt f\"ur die
\"ubergehende Ionenspezies $i$:
\[
E_{\mathrm{Donnan}} = \phi^\alpha - \phi^\beta
= \frac{RT}{z_i F} \ln \frac{a_i^\beta}{a_i^\alpha}.
\]
\end{lemma}

\begin{proof}
Im Gleichgewicht: $\tilde{\mu}_i^\alpha = \tilde{\mu}_i^\beta$.
Also:
\begin{align*}
\mu_i^{0,\alpha} + RT\ln a_i^\alpha + z_i F \phi^\alpha
&= \mu_i^{0,\beta} + RT\ln a_i^\beta + z_i F \phi^\beta.
\end{align*}
F\"ur die gleiche Ionenspezies in beiden Phasen gilt $\mu_i^{0,\alpha} = \mu_i^{0,\beta}$
(selbe Standardbedingungen im Wassermedium).
Umstellen:
\[
z_i F (\phi^\alpha - \phi^\beta) = RT\ln a_i^\beta - RT\ln a_i^\alpha
= RT\ln\frac{a_i^\beta}{a_i^\alpha}.
\]
Division durch $z_i F$:
\[
E_{\mathrm{Donnan}} = \phi^\alpha - \phi^\beta = \frac{RT}{z_i F}\ln\frac{a_i^\beta}{a_i^\alpha}. \qquad \blacksquare
\]
\end{proof}

\begin{bemerkung}
Das Donnan-Potential ist der Schlu\"ssselmechanismus f\"ur das
ionotronische NOT-Gatter (Satz~\ref{satz:kohlrausch}):
Ein hoher Konzentrationsunterschied erzeugt eine Sperrspannung,
die den Kanalstrom blockiert.
\end{bemerkung}

%%%%%%%%%%%%%%%%%%%%%%%%%%%%%%%%%%%
\section{Die Ionotronische Turing-Maschine}
%%%%%%%%%%%%%%%%%%%%%%%%%%%%%%%%%%%

\subsection{Formale Definition}

\begin{definition}[Ionotronische Turing-Maschine]
\label{def:itm}
Eine \emph{Ionotronische Turing-Maschine} (ITM) ist ein 7-Tupel
\[
\ITM = (Q,\, \Sigma,\, \Gamma,\, \delta,\, q_0,\, q_{\mathrm{halt}},\, \cth),
\]
wobei:
\begin{itemize}
    \item $Q = \{q_0, q_1, \ldots, q_{|Q|-1}\}$: endliche Menge von
          \emph{Zust\"anden}, realisiert als diskrete Elektrodenkonfigurationen,
    \item $\Sigma \subsetneq \Gamma$: \emph{Eingabealphabet},
    \item $\Gamma = \{C_0, C_1, \ldots, C_{|\Gamma|-2}, B\}$: \emph{Bandalphabet}
          mit $|\Gamma|$ diskreten Ionenkonzentrationsniveaus und
          Blanksymbol $B$ (destilliertes Wasser, $c = 0$),
    \item $\delta: (Q \setminus \{q_{\mathrm{halt}}\}) \times \Gamma \to
          Q \times \Gamma \times \{L, R, S\}$: \emph{\"Ubergangsrelation}
          (partiell),
    \item $q_0 \in Q$: \emph{Startzustand},
    \item $q_{\mathrm{halt}} \in Q$: \emph{Haltezustand} ($q_{\mathrm{halt}} \neq q_0$),
    \item $\cth > 0$: \emph{Konzentrationsschwelle}, definiert die Bit-Grenze:
          $c < \cth \Leftrightarrow$ Symbol~$= C_0$; $c \geq \cth \Leftrightarrow$
          Symbol~$= C_1$ (im bin\"aren Fall).
\end{itemize}
\end{definition}

\begin{definition}[Ionisches Band]
\label{def:band}
Das \emph{ionische Band} der ITM ist ein bidirektional unendliches
Mikrokanalarray:
\[
\mathcal{B} = \{b_i\}_{i \in \mathbb{Z}}, \quad b_i \in \Gamma.
\]
Die physikalische Realisierung: Zelle $i$ ist ein Mikrokanalabschnitt
der L\"ange $\ell_{\mathrm{cell}} = 50\,\mu$m, Breite $w = 10\,\mu$m,
H\"ohe $h = 5\,\mu$m, gef\"ullt mit Elektrolyt der Konzentration
$c_i \in \{0, \cth\}$ (bin\"arer Fall).
\end{definition}

\begin{definition}[Konfiguration der ITM]
\label{def:config}
Eine \emph{Konfiguration} der ITM ist ein Tripel
\[
\mathrm{Konf} = (q,\, \mathcal{B},\, k) \in Q \times \Gamma^{\mathbb{Z}} \times \mathbb{Z},
\]
mit Zustand $q$, Bandinhalt $\mathcal{B}$ und Kopfposition $k \in \mathbb{Z}$.
\end{definition}

\begin{definition}[Berechnungsschritt]
\label{def:schritt}
Sei $\delta(q, b_k) = (q', b_k', D)$.
Der \emph{Berechnungsschritt} $\vdash$ ist:
\[
(q, \mathcal{B}, k) \;\vdash\; (q', \mathcal{B}', k'),
\]
wobei:
\[
\mathcal{B}'(i) = \begin{cases} b_k' & i = k, \\ b_i & i \neq k, \end{cases}
\qquad
k' = \begin{cases} k+1 & D = R, \\ k-1 & D = L, \\ k & D = S. \end{cases}
\]
$\vdash^*$ bezeichnet den reflexiv-transitiven Abschluss von $\vdash$.
\end{definition}

\begin{definition}[Berechnung und Akzeptanz]
\label{def:berechnung}
Die ITM \emph{berechnet} $f: \{0,1\}^* \to \{0,1\}^*$, wenn sie f\"ur
jede Eingabe $x \in \{0,1\}^*$ (kodiert als Anfangsbandinhalt) in endlich
vielen Schritten in den Zustand $q_{\mathrm{halt}}$ \"ubergeht und das Band
$f(x)$ enth\"alt.
\end{definition}

\begin{figure}[h!]
\centering
\includegraphics[width=0.95\textwidth]{plot2_itm_band.pdf}
\caption{Links: Zustandsgraph der ITM mit 5 Zust\"anden.
Rechts: Ionisches Band als Mikrokanalarray; Konzentrationsniveaus
kodieren Bandsymbole.}
\label{fig:itm_band}
\end{figure}

\subsection{Kodierungslemma}

\begin{lemma}[Injektivit\"at der ionischen Kodierung]
\label{lem:inject}
Die Kodierungsabbildung $\mathrm{enc}: \{0,1\} \to \{c_{\mathrm{low}}, c_{\mathrm{high}}\}$,
\[
\mathrm{enc}(0) = c_{\mathrm{low}} < \cth, \quad \mathrm{enc}(1) = c_{\mathrm{high}} > \cth,
\]
und ihre Inverse $\mathrm{dec}(c) = \mathbf{1}[c \geq \cth]$
sind wohldefiniert und zueinander invers, sofern
$c_{\mathrm{low}} < \cth < c_{\mathrm{high}}$.

Dar\"uber hinaus ist die Kodierung \emph{rauschstabil}: F\"ur thermisches
Konzentrationsrauschen $\varepsilon \sim \mathcal{N}(0, \sigma_c^2)$ mit
$\sigma_c \ll (\cth - c_{\mathrm{low}})$ und
$\sigma_c \ll (c_{\mathrm{high}} - \cth)$ gilt:
\[
P[\mathrm{Dekodierungsfehler}] = P[c + \varepsilon \text{ falsch dekodiert}]
\leq 2\,\exp\!\left(-\frac{(\cth - c_{\mathrm{low}})^2}{2\sigma_c^2}\right).
\]
\end{lemma}

\begin{proof}
\textbf{Wohldefinierte Kodierung:}
Da $c_{\mathrm{low}} < \cth < c_{\mathrm{high}}$ (Voraussetzung),
gilt $\mathrm{dec}(\mathrm{enc}(0)) = \mathbf{1}[c_{\mathrm{low}} \geq \cth] = 0$
und $\mathrm{dec}(\mathrm{enc}(1)) = \mathbf{1}[c_{\mathrm{high}} \geq \cth] = 1$.
Injektivit\"at: $\mathrm{enc}(0) = c_{\mathrm{low}} \neq c_{\mathrm{high}} = \mathrm{enc}(1)$.

\textbf{Rauschstabilit\"at:}
Fehler bei Kodierung~0 tritt auf, wenn $c_{\mathrm{low}} + \varepsilon \geq \cth$,
d.\,h.\ $\varepsilon \geq \cth - c_{\mathrm{low}} =: \Delta_0 > 0$.
\[
P[\varepsilon \geq \Delta_0] = P\!\left[\frac{\varepsilon}{\sigma_c} \geq \frac{\Delta_0}{\sigma_c}\right]
\leq \exp\!\left(-\frac{\Delta_0^2}{2\sigma_c^2}\right)
\]
nach der Gaussian-Tail-Schranke ($P[Z \geq t] \leq e^{-t^2/2}$ f\"ur $Z \sim \mathcal{N}(0,1)$).
Analog: Fehler bei Kodierung~1: $\varepsilon \leq \cth - c_{\mathrm{high}} < 0$,
Wahrscheinlichkeit $\leq \exp(-\Delta_1^2/(2\sigma_c^2))$ mit $\Delta_1 = c_{\mathrm{high}} - \cth$.
Union Bound:
\[
P[\text{Fehler}] \leq \exp\!\left(-\frac{\Delta_0^2}{2\sigma_c^2}\right)
+ \exp\!\left(-\frac{\Delta_1^2}{2\sigma_c^2}\right)
\leq 2\exp\!\left(-\frac{\min(\Delta_0,\Delta_1)^2}{2\sigma_c^2}\right). \qquad \blacksquare
\]
\end{proof}

%%%%%%%%%%%%%%%%%%%%%%%%%%%%%%%%%%%
\section{Ionische Logikgatter}
%%%%%%%%%%%%%%%%%%%%%%%%%%%%%%%%%%%

\subsection{Kohlrausch-Schaltungssatz}

\begin{satz}[Kohlrausch-Schaltungssatz]
\label{satz:kohlrausch}
In einem ionotronischen Schaltkreis mit steuerbaren Kanalkonzentrationen
lassen sich die Booleschen Logikfunktionen AND, OR und NOT physikalisch
realisieren:

\begin{enumerate}[label=(\roman*)]
\item \textbf{AND-Gatter (Serienschaltung):}
      F\"ur zwei serielle Kan\"ale $K_1$ und $K_2$ mit Leitf\"ahigkeiten
      $\kappa_1 = \kappa(\mathrm{enc}(A))$ und $\kappa_2 = \kappa(\mathrm{enc}(B))$
      gilt mit Schwelle $\kth = \kappa(\chigh)/4$:
      \[
      \mathrm{dec}(\kges) = A \wedge B,
      \]
      wobei $\kges = (\kappa_1^{-1} + \kappa_2^{-1})^{-1}$ die
      Serien-Leitf\"ahigkeit.

\item \textbf{OR-Gatter (Parallelschaltung):}
      F\"ur zwei parallele Kan\"ale gilt mit $\kth = \kappa(\chigh)/2$:
      \[
      \mathrm{dec}(\kges) = A \vee B,
      \]
      wobei $\kges = \kappa_1 + \kappa_2$ die Parallel-Leitf\"ahigkeit.

\item \textbf{NOT-Gatter (Donnan-Sperrung):}
      F\"ur Eingang $A$ und feste Gate-Konzentration $c_{\mathrm{gate}} = \chigh$:
      \[
      \mathrm{dec}(\kappa_{\mathrm{out}}) = \neg A.
      \]
\end{enumerate}
\end{satz}

\begin{proof}
Wir bezeichnen $\kappa_0 := \kappa(\clow) \approx 0$ und
$\kappa_1 := \kappa(\chigh) > 0$.

\textbf{Beweis (i) AND -- Serienschaltung:}

Die vier Eingangsf\"alle:
\begin{itemize}
\item $(A,B) = (0,0)$: $\kappa_1 = \kappa_0 \approx 0$, $\kappa_2 = \kappa_0 \approx 0$.
  $\kges = (\kappa_0^{-1}+\kappa_0^{-1})^{-1} \approx 0 \ll \kth$. Ausgang: 0. \checkmark
\item $(A,B) = (1,0)$: $\kappa_1 = \kappa_1$, $\kappa_2 = \kappa_0 \approx 0$.
  $\kges = (\kappa_1^{-1} + \kappa_0^{-1})^{-1} \approx 0 \ll \kth$
  (da $\kappa_0^{-1} \to \infty$ dominiert). Ausgang: 0. \checkmark
\item $(A,B) = (0,1)$: Symmetrisch zu $(1,0)$. Ausgang: 0. \checkmark
\item $(A,B) = (1,1)$: $\kappa_1 = \kappa_2 = \kappa_1$.
  $\kges = (\kappa_1^{-1}+\kappa_1^{-1})^{-1} = \kappa_1/2 > \kappa_1/4 = \kth$.
  Ausgang: 1. \checkmark
\end{itemize}
Alle vier F\"alle stimmen mit der AND-Wahrheitstabelle \"uberein.

\textbf{Beweis (ii) OR -- Parallelschaltung:}

\begin{itemize}
\item $(A,B) = (0,0)$: $\kges = \kappa_0 + \kappa_0 = 2\kappa_0 \approx 0 < \kth$. Ausgang: 0. \checkmark
\item $(A,B) = (1,0)$: $\kges = \kappa_1 + \kappa_0 \approx \kappa_1 > \kappa_1/2 = \kth$. Ausgang: 1. \checkmark
\item $(A,B) = (0,1)$: Symmetrisch. Ausgang: 1. \checkmark
\item $(A,B) = (1,1)$: $\kges = 2\kappa_1 > \kth$. Ausgang: 1. \checkmark
\end{itemize}

\textbf{Beweis (iii) NOT -- Donnan-Sperrung:}

Das NOT-Gatter nutzt das Donnan-Gleichgewicht (Lemma~\ref{lem:donnan}).
Sei $c_{\mathrm{gate}} = \chigh$ die fest eingestellte Tor-Konzentration.
\begin{itemize}
\item Eingang $A = 0$ ($c_{\mathrm{in}} = \clow$):
  Donnan-Spannung: $E_D = \frac{RT}{F}\ln\frac{\chigh}{\clow} > 0$.
  Da $E_D > 0$, wird der Ausgangskanal von Ionen \emph{bef\"ullt}
  (positives Donnan-Potential treibt Ionen in den Kanal).
  Also $c_{\mathrm{out}} = \chigh \geq \cth$: Ausgang 1. \checkmark
\item Eingang $A = 1$ ($c_{\mathrm{in}} = \chigh$):
  $E_D = \frac{RT}{F}\ln\frac{\chigh}{\chigh} = 0$.
  Kein Donnan-Antrieb; Osmotischer Druck gleicht sich aus.
  Da zus\"atzlich der Gate-Gegenstrom bei gleichen Konzentrationen
  keine Nettoionen in den Ausgang treibt (elektrochemisches Gleichgewicht),
  bleibt $c_{\mathrm{out}} = \clow < \cth$: Ausgang 0. \checkmark
\end{itemize}
Damit ist NOT realisiert. \checkmark
\end{proof}

\begin{figure}[h!]
\centering
\includegraphics[width=0.95\textwidth]{plot3_logikgatter.pdf}
\caption{Wahrheitstabellen der sechs ionischen Grundgatter.
Farbige Zellen: Ausgang~1; helle Zellen: Ausgang~0.}
\label{fig:logikgatter}
\end{figure}

\begin{satz}[Funktionale Vollst\"andigkeit ionischer Gatter]
\label{satz:vollstaendig}
Das Gatter-System $\mathcal{G} = \{\mathrm{AND}, \mathrm{NOT}\}$ ionischer
Logikgatter (Satz~\ref{satz:kohlrausch}) ist \emph{funktional vollst\"andig}:
Jede Boole'sche Funktion $f: \{0,1\}^n \to \{0,1\}$ ist durch eine endliche
Schaltung aus AND- und NOT-Gattern realisierbar.
\end{satz}

\begin{proof}
\textbf{Schritt 1: OR aus AND und NOT.}
Nach den de~Morgan-Gesetzen:
\[
A \vee B = \neg(\neg A \wedge \neg B).
\]
Die rechte Seite verwendet nur NOT und AND, ist also in $\mathcal{G}$ realisierbar.

\textbf{Schritt 2: Disjunktive Normalform (DNF).}
Jede Boole'sche Funktion $f:\{0,1\}^n \to \{0,1\}$ l\"asst sich als DNF schreiben:
\[
f(x_1,\ldots,x_n) = \bigvee_{x \in f^{-1}(1)} \left(\bigwedge_{i=1}^n \ell_i(x)\right),
\]
wobei $\ell_i(x) = x_i$ falls $x_i = 1$ und $\ell_i(x) = \neg x_i$ falls $x_i = 0$
(jeder Minterm).

\textbf{Schritt 3: Realisierung in $\mathcal{G}$.}
Aus Schritt~1 ist OR durch AND und NOT darstellbar.
Jeder Minterm $\bigwedge_{i=1}^n \ell_i$ verwendet AND und (potenziell) NOT.
Die \"au{\ss}ere Disjunktion der Minterme wird durch OR = AND+NOT realisiert.
Damit ist $f$ in $\mathcal{G}$ realisierbar.

\textbf{Schritt 4: Endlichkeit der Schaltung.}
Die Anzahl der Minterme ist $|f^{-1}(1)| \leq 2^n$.
Jeder Minterm ben\"otigt $\leq n$ NOT-Gatter und $\leq n-1$ AND-Gatter.
Die Disjunktion von $m$ Mintermen ben\"otigt $\leq m-1$ OR-Gatter,
jedes darstellbar durch 2 NOT + 1 AND.
Gesamte Gatterzahl: $\leq 2^n \cdot (2n-1) + 3 \cdot (2^n-1) < \infty$.
\end{proof}

\begin{korollar}[NAND-Vollst\"andigkeit]
\label{kor:nand}
Das einzelne NAND-Gatter $\{(A,B) \mapsto \neg(A \wedge B)\}$ ist allein
funktional vollst\"andig.
\end{korollar}

\begin{proof}
NOT aus NAND: $\neg A = A \;\mathrm{NAND}\; A$ (da $(A \wedge A) = A$).
AND aus NAND: $A \wedge B = \neg(A \;\mathrm{NAND}\; B) = (A\;\mathrm{NAND}\;B)\;\mathrm{NAND}\;(A\;\mathrm{NAND}\;B)$.
Da $\{\mathrm{NOT}, \mathrm{AND}\}$ vollst\"andig (Satz~\ref{satz:vollstaendig}),
und beide durch NAND darstellbar, ist NAND allein vollst\"andig.
\end{proof}

%%%%%%%%%%%%%%%%%%%%%%%%%%%%%%%%%%%
\section{Physikalische Realisierbarkeit}
%%%%%%%%%%%%%%%%%%%%%%%%%%%%%%%%%%%

\begin{satz}[Physikalische Realisierbarkeit]
\label{satz:phys_real}
Jeder Berechnungsschritt $\delta(q, b) = (q', b', D)$ der ITM ist
physikalisch in einem mikrofluidischen System mit den folgenden
Operationen realisierbar:

\begin{enumerate}[label=(\roman*)]
\item \textbf{Lesen:} Impedanzspektroskopie liefert $\kappa_i$ in $O(1)$ Zeit;
      $b = \mathrm{dec}(\kappa_i) \in \{0,1\}$ mit Fehlerrate
      $P_{\mathrm{err}} \leq 2\exp(-\Delta_0^2/(2\sigma_c^2))$ (Lemma~\ref{lem:inject}).
\item \textbf{Schreiben:} Mikropumpen-Injektion von $\Delta V = (\mathrm{enc}(b') - c_i) \cdot V_{\mathrm{cell}}$
      Liter NaCl-L\"osung in Zeitkonstante $\tau_{\mathrm{write}} = V_{\mathrm{cell}}/Q_{\mathrm{pump}}$.
\item \textbf{Kopfbewegung:} Elektrophoretische Verschiebung der Elektroden
      um $w = 10\,\mu$m in Zeit $\tau_{\mathrm{move}} = w/v_{\mathrm{EP}}$
      mit elektrophoretischer Geschwindigkeit $v_{\mathrm{EP}} = \mu_e E$.
\item \textbf{Zustands\"anderung:} Umschalten der Gate-Elektrodenspannung
      in Zeit $\tau_{\mathrm{state}} \ll \tau_{\mathrm{write}}$ (elektronische Schaltzeit).
\end{enumerate}
\end{satz}

\begin{proof}
\textbf{Zu (i) -- Lesen:}
Die Impedanzspektroskopie misst den komplexen Widerstand $Z(\omega) = R + j/(\omega C)$
des Kanals. Im Niederfrequenzlimit $\omega \to 0$: $Z \to R = (\kappa A/L)^{-1}$
(Widerstand des Kanals der Querschnittsfl\"ache $A$ und L\"ange $L$).
Daher $\kappa = L/(R \cdot A)$; Messung von $R$ liefert $\kappa$ in $O(1)$.
Fehlerrate: folgt aus Lemma~\ref{lem:inject}.

\textbf{Zu (ii) -- Schreiben:}
Das Zellvolumen $V_{\mathrm{cell}} = w \cdot h \cdot \ell_{\mathrm{cell}} = 10 \cdot 5 \cdot 50\,\mu$m$^3
= 2500\,\mu$m$^3 = 2{,}5 \cdot 10^{-15}$\,L.
F\"ur eine Mikropumpe mit Durchfluss $Q_{\mathrm{pump}} = 1\,\text{nL/s} = 10^{-9}\,\text{L/s}$
(realistischer Wert~\cite{laser2004}):
\[
\tau_{\mathrm{write}} = \frac{V_{\mathrm{cell}}}{Q_{\mathrm{pump}}}
= \frac{2{,}5 \cdot 10^{-15}}{10^{-9}}\,\text{s} = 2{,}5\,\mu\text{s}.
\]
Dies ist physikalisch realisierbar.

\textbf{Zu (iii) -- Kopfbewegung:}
Die elektrophoretische Mobilit\"at von Na$^+$:
$\mu_e = 5{,}19 \cdot 10^{-8}\,\text{m}^2/(\text{V\,s})$.
Bei $E = 10^4\,\text{V/m}$ (1\,V \"uber 100\,$\mu$m):
$v_{\mathrm{EP}} = \mu_e E = 5{,}19 \cdot 10^{-4}\,\text{m/s}$.
Bewegung um $w = 10^{-5}\,\text{m}$:
\[
\tau_{\mathrm{move}} = \frac{w}{v_{\mathrm{EP}}} = \frac{10^{-5}}{5{,}19 \cdot 10^{-4}}\,\text{s}
\approx 19{,}3\,\text{ms}.
\]
Dies ist der Flaschenhals des Systems; f\"ur praktische Anwendungen
sind optische Aktuierungsmethoden vorzuziehen (Faktor $10^3$ schneller).

\textbf{Zu (iv) -- Zustands\"anderung:}
Elektronische Gate-Umschaltung erfolgt auf Nanosekunden-Skala
($\tau_{\mathrm{state}} \approx 1\,\text{ns}$), vernachl\"assigbar gegen\"uber
$\tau_{\mathrm{write}}$.

Da alle vier Operationen physikalisch realisierbar sind, ist jeder
Berechnungsschritt realisierbar.
\end{proof}

\begin{korollar}[Zeitkomplexiti\"at der ITM]
\label{kor:zeitkompl}
Die physikalische Ausf\"uhrungszeit f\"ur $T$ Berechnungsschritte der ITM ist:
\[
t_{\mathrm{total}} = T \cdot (\tau_{\mathrm{read}} + \tau_{\mathrm{write}} + \tau_{\mathrm{move}})
= T \cdot O(\tau_{\mathrm{move}}) = O(T \cdot w / v_{\mathrm{EP}}).
\]
F\"ur eine optimierte Implementierung mit $v_{\mathrm{EP}} = 10^{-3}\,\text{m/s}$:
$\tau_{\mathrm{move}} \approx 10\,\text{ms}$ pro Schritt.
\end{korollar}

%%%%%%%%%%%%%%%%%%%%%%%%%%%%%%%%%%%
\section{Das Analog-Paradigma und Shannon-Theorie}
%%%%%%%%%%%%%%%%%%%%%%%%%%%%%%%%%%%

\subsection{Analoger Informationsgehalt ionischer Kan\"ale}

\begin{satz}[Analoger Informationsgehalt]
\label{satz:analog_info}
Ein ionischer Kanal der L\"ange $\ell$ mit Konzentrationsprofil
$c(x) \in [0, c_{\max}]$ und thermischem Rauschen
$\sigma_c^2 = k_B T \cdot c / (z^2 F^2 / (R^2 T^2))$ tr\"agt
einen analogen Informationsgehalt:
\[
I_{\mathrm{analog}} = \log_2\!\left(\frac{c_{\max}}{c_{\mathrm{noise}}}\right) \cdot \frac{\ell}{\ell_{\mathrm{cell}}}
= \log_2(\mathrm{DR}) \cdot \frac{\ell}{\ell_{\mathrm{cell}}} \quad [\mathrm{bit}],
\]
wobei $\mathrm{DR} = c_{\max}/c_{\mathrm{noise}}$ der \emph{Dynamikbereich} ist.
\end{satz}

\begin{proof}
\textbf{Schritt 1: Rauschmodell.}
Thermische Fluktuationen der Ionenkonzentration in einem Volumen $V$:
Nach Einstein-Smoluchowski gilt f\"ur die mittlere quadratische Konzentrationsfluktuation:
\[
\langle (\delta c)^2 \rangle = \frac{c}{N_A V},
\]
wobei $N_A$ die Avogadro-Zahl.
F\"ur $V = V_{\mathrm{cell}} = 2{,}5 \cdot 10^{-15}\,\text{L}$ und
$c = 0{,}1\,\text{mol/L}$:
\[
\sigma_c = \sqrt{\frac{c}{N_A V}} = \sqrt{\frac{0{,}1}{6{,}022 \cdot 10^{23} \cdot 2{,}5 \cdot 10^{-18}}}
\approx 2{,}6 \cdot 10^{-4}\,\text{mol/L}.
\]

\textbf{Schritt 2: Informationskapazit\"at.}
Modelliere den ionischen Kanal als Gauß'schen Kanal mit Signal $c \in [0, c_{\max}]$
und Rauschen $\varepsilon \sim \mathcal{N}(0, \sigma_c^2)$.
Nach Shannon~\cite{shannon1948}: Kanalkapazit\"at
$C = \frac{1}{2}\log_2(1 + c_{\max}^2/\sigma_c^2)$.
F\"ur $c_{\max} \gg \sigma_c$: $C \approx \log_2(c_{\max}/\sigma_c) = \log_2(\mathrm{DR})$.

\textbf{Schritt 3: Skalierung \"uber L\"ange.}
Aufeinanderfolgende Zellen \"ubertragen unabh\"angige Information
(bei ausreichendem Abstand $\ell_{\mathrm{cell}}$).
F\"ur $N = \ell/\ell_{\mathrm{cell}}$ Zellen:
$I_{\mathrm{analog}} = N \cdot C \approx (\ell/\ell_{\mathrm{cell}}) \cdot \log_2(\mathrm{DR})$.
\end{proof}

\begin{satz}[Diskretisierungsverlust]
\label{satz:diskr_verlust}
Bei der $n$-bit-Digitalisierung eines ionischen Signals mit Dynamikbereich
$\mathrm{DR} = 2^m$ ($m > n$ bit \"aquivalent) entsteht ein
Signal-zu-Quantisierungsrauschen-Verh\"altnis:
\[
\mathrm{SQNR}_n = 6{,}02\,n + 1{,}76\,[\mathrm{dB}],
\]
und der \emph{Informationsverlust} pro Zelle:
\[
\Delta I_{\mathrm{cell}} = \log_2(\mathrm{DR}) - n = m - n\,[\mathrm{bit}].
\]
F\"ur $n < m$ geht Information \emph{unwiederbringlich} verloren.
\end{satz}

\begin{proof}
\textbf{SQNR-Formel:}
Ein $n$-bit-Quantisierer mit Vollaussteuerung hat Quantisierungsrauschen
$P_q = \Delta^2/12 = (c_{\max}/(2^n))^2/12$,
wobei $\Delta = c_{\max}/2^n$ die Stufenbreite.
Signalleistung f\"ur Vollaussteuerung mit Sinussignal: $P_s = c_{\max}^2/8$.
\[
\mathrm{SQNR} = \frac{P_s}{P_q} = \frac{c_{\max}^2/8}{c_{\max}^2/(12 \cdot 2^{2n})}
= \frac{12 \cdot 2^{2n}}{8} = \frac{3}{2} \cdot 4^n.
\]
In Dezibel: $\mathrm{SQNR}_n = 10\log_{10}(1{,}5 \cdot 4^n)
= 10\log_{10}(1{,}5) + 20n\log_{10}(2) \approx 1{,}76 + 6{,}02\,n\,\text{dB}$.

\textbf{Informationsverlust:}
Der Kanal hat analoge Kapazit\"at $m$ bit (\"aquivalent), der Quantisierer
liefert $n$ bit. Verlust: $m - n > 0$ bei $n < m$.
Der Verlust ist irreversibel (Datenkompression mit Verlust ist injektiv
auf dem quantisierten Gitter, aber surjektiv auf das analoge Intervall
ist ohne zus\"atzliche Information nicht m\"oglich).
\end{proof}

\begin{figure}[h!]
\centering
\includegraphics[width=0.95\textwidth]{plot4_analog_digital.pdf}
\caption{Analog vs. Digital in der Ionotronik.
Oben links: Nernst-Potential als stetige Funktion.
Oben rechts: Quantisierungsfehler verschiedener Bittiefen.
Unten links: Shannon-Kanalkapazit\"at ionischer und CMOS-Kan\"ale.
Unten rechts: Energieverbrauch-Vergleich.}
\label{fig:analog_digital}
\end{figure}

%%%%%%%%%%%%%%%%%%%%%%%%%%%%%%%%%%%
\section{Ionisch-photonische Hybridarchitektur}
%%%%%%%%%%%%%%%%%%%%%%%%%%%%%%%%%%%

\subsection{Plasmonische Grundlagen}

\begin{definition}[Lokalisierte Oberfl\"achenplasmonresonanz]
\label{def:lspr}
Die \emph{lokalisierte Oberfl\"achenplasmonresonanz} (LSPR) metallischer
Nanopartikel (Radius $R \ll \lambda$) liegt bei der Wellenl\"ange $\lambda_{\mathrm{res}}$,
die die Mie-Resonanzbedingung erf\"ullt:
\[
\varepsilon_{\mathrm{Metall}}(\lambda_{\mathrm{res}}) = -2\,\varepsilon_{\mathrm{Medium}},
\]
wobei $\varepsilon_{\mathrm{Metall}} = \varepsilon_1 + i\varepsilon_2$ die
(komplexe) Dielektrizit\"atskonstante des Metalls und
$\varepsilon_{\mathrm{Medium}} = n_{\mathrm{eff}}^2$ die des umgebenden Mediums.
\end{definition}

\begin{lemma}[Linearit\"at der Resonanzverschiebung]
\label{lem:linear_shift}
F\"ur kleine \"Anderungen $\Delta n = n_{\mathrm{eff}} - n_0 \ll n_0$ gilt
die Linearit\"at der Resonanzverschiebung:
\[
\Delta\lambda_{\mathrm{res}} = S \cdot \Delta n + O(\Delta n^2),
\]
wobei $S = d\lambda_{\mathrm{res}}/dn|_{n=n_0}$ die \emph{refraktometrische Sensitivit\"at}
[nm/RIU] (RIU = Refractive Index Unit).
\end{lemma}

\begin{proof}
Differenziere die Mie-Resonanzbedingung $\varepsilon_1(\lambda) = -2n^2$ nach $n$:
\[
\frac{d\varepsilon_1}{d\lambda}\bigg|_{\lambda_{\mathrm{res}}} \cdot \frac{d\lambda_{\mathrm{res}}}{dn}
= -4n.
\]
Aufl\"osen:
\[
S = \frac{d\lambda_{\mathrm{res}}}{dn} = \frac{-4n}{d\varepsilon_1/d\lambda\big|_{\lambda_{\mathrm{res}}}}.
\]
F\"ur Gold-Nanopartikel bei $\lambda_{\mathrm{res}} \approx 520\,\text{nm}$:
$d\varepsilon_1/d\lambda \approx -0{,}5\,\text{nm}^{-1}$, $n_0 = 1{,}33$:
$S \approx 4 \cdot 1{,}33/0{,}5 \approx 120\,\text{nm/RIU}$.
Taylor-Entwicklung liefert den $O(\Delta n^2)$-Fehlerterm.
\end{proof}

\begin{satz}[Ionisch-photonischer Hybridisierungssatz]
\label{satz:hybrid}
In einem ionoxoqtum-System mit ionischer Brechungsindex-\"Anderung
$\Delta n(c) = (dn/dc) \cdot c$ und Sensitivit\"at $S$ gilt:
\[
\Delta\lambda_{\mathrm{res}}(c) = S \cdot \frac{dn}{dc} \cdot c + O(c^2).
\]
F\"ur NaCl mit $dn/dc = 0{,}002\,\text{L/mol}$ und Gold-Nanopartikeln
($S = 120\,\text{nm/RIU}$) ist der Linearit\"atsbereich $c < 0{,}5\,\text{mol/L}$
mit Fehler $< 1\%$.
Damit ist das ionische Bit durch Spektroskopie kontaktlos auslesbar,
wenn $\Delta\lambda_{\mathrm{res}}(c_{\mathrm{high}}) > \Delta\lambda_{\min} = 0{,}1\,\text{nm}$,
d.\,h.\ $c_{\mathrm{high}} > 0{,}42\,\text{mol/L}$.
\end{satz}

\begin{proof}
\textbf{Brechungsindex-\"Anderung:}
F\"ur NaCl-L\"osungen gilt bis $c = 1\,\text{mol/L}$~\cite{quan2009}:
$n(c) = n_0 + (dn/dc)\cdot c + O(c^2)$ mit $n_0 = 1{,}3330$ und $dn/dc = 0{,}0018$--$0{,}0022\,\text{L/mol}$.
Wir nehmen $dn/dc = 0{,}002\,\text{L/mol}$.

\textbf{Resonanzverschiebung:}
Aus Lemma~\ref{lem:linear_shift}:
\[
\Delta\lambda_{\mathrm{res}}(c) = S \cdot \Delta n(c) = S \cdot \frac{dn}{dc} \cdot c
= 120 \cdot 0{,}002 \cdot c = 0{,}24\,c\,[\text{nm}],
\]
f\"ur $c$ in mol/L.

\textbf{Aufl\"osungsbedingung:}
$\Delta\lambda_{\mathrm{res}} > \Delta\lambda_{\min} = 0{,}1\,\text{nm}$
$\Leftrightarrow 0{,}24\,c > 0{,}1 \Leftrightarrow c > 0{,}42\,\text{mol/L}$.

\textbf{Fehlerterm:}
Der $O(c^2)$-Term aus der Taylor-Entwicklung von $n(c)$ ist
$\leq (d^2n/dc^2)\cdot c^2/2 \approx 10^{-4}\cdot c^2\,\text{L}^2/\text{mol}^2$.
F\"ur $c \leq 0{,}5\,\text{mol/L}$: Fehler $\leq 2{,}5 \cdot 10^{-5} \ll 0{,}002\cdot c$:
relativer Fehler $< 1\%$.
\end{proof}

\begin{figure}[h!]
\centering
\includegraphics[width=0.95\textwidth]{plot5_ionoxoqtum.pdf}
\caption{Links: Plasmonen-Transmission f\"ur verschiedene NaCl-Konzentrationen.
Rechts: F\"unfschichtige ionoxoqtum-Architektur.}
\label{fig:ionoxoqtum}
\end{figure}

%%%%%%%%%%%%%%%%%%%%%%%%%%%%%%%%%%%
\section{Turing-Vollst\"andigkeit der ITM}
%%%%%%%%%%%%%%%%%%%%%%%%%%%%%%%%%%%

\subsection{\"Aquivalenz von ITM und klassischer TM}

\begin{satz}[Ionotronischer Universalit\"atssatz]
\label{satz:itm_universal}
Sei $\mathcal{T} = (Q_T, \Sigma_T, \Gamma_T, \delta_T, q_0^T, q_{\mathrm{halt}}^T)$
eine beliebige (klassische) Turing-Maschine.
Dann existiert eine Ionotronische Turing-Maschine $\ITM$ mit:
\begin{enumerate}[label=(\roman*)]
    \item $\ITM$ simuliert $\mathcal{T}$ Schritt f\"ur Schritt,
    \item $\ITM$ h\"alt auf Eingabe $x$ gdw.\ $\mathcal{T}$ h\"alt auf $x$,
    \item Ausf\"uhrungszeit: $t_{\mathrm{ITM}}(n) = O(t_{\mathcal{T}}(n))$
          (linearer \"Overhead, da jeder TM-Schritt in $O(1)$ ITM-Schritten
          simuliert wird).
\end{enumerate}
Insbesondere: ITM $\equiv_T$ TM (Turing-\"Aquivalenz).
\end{satz}

\begin{proof}
Wir konstruieren $\ITM$ explizit als Simulation von $\mathcal{T}$.

\textbf{Konstruktion 1: Zustandsmenge.}
Setze $Q := Q_T \cup \{q_{\mathrm{halt}}^{\mathrm{ITM}}\}$, wobei
$q_{\mathrm{halt}}^{\mathrm{ITM}}$ eine Kopie von $q_{\mathrm{halt}}^T$ ist.
Die Zust\"ande von $\ITM$ entsprechen denen von $\mathcal{T}$;
physikalisch: jeder Zustand $q \in Q_T$ wird durch eine eindeutige
Elektroden-Spannungs-Konfiguration $\mathbf{V}_q \in \mathbb{R}^{|Q_T|}$
repr\"asentiert.

\textbf{Konstruktion 2: Bandalphabet-Kodierung.}
Definiere $\Gamma := \Gamma_T$ (gleiche Symbolmengen).
Die Kodierung: W\"ahle $\cth$ und diskrete Konzentrationsniveaus
$\{C_s \mid s \in \Gamma_T\}$ so, dass
$|C_{s_1} - C_{s_2}| \geq 2\sigma_c$ f\"ur alle $s_1 \neq s_2$
(Trennbarkeit), d.\,h.\
\[
C_s := s \cdot \frac{c_{\max}}{|\Gamma_T| - 1}, \quad s = 0, 1, \ldots, |\Gamma_T|-1,
\]
mit $c_{\max} = (|\Gamma_T|-1) \cdot 2\sigma_c + \cth$ (ausreichend gro{\ss}).

\textbf{Konstruktion 3: \"Ubergangsfunktion.}
Setze $\delta_{\mathrm{ITM}}(q, C_s) := (q', C_{s'}, D)$
f\"ur $(q', s', D) = \delta_T(q, s)$.
Dies ist wohldefiniert, da $\delta_T$ auf $Q_T \times \Gamma_T$ definiert
und die Kodierung bijektiv ist.

\textbf{Beweis der Simulation (Schritt-f\"ur-Schritt):}
Behauptung: F\"ur jede Konfiguration $(q, \mathcal{B}_T, k)$ von $\mathcal{T}$
gilt: Die entsprechende Konfiguration $(q, \mathcal{B}_{\mathrm{ITM}}, k)$ von $\ITM$
(mit $\mathcal{B}_{\mathrm{ITM}}(i) = C_{\mathcal{B}_T(i)}$ f\"ur alle $i$)
vollf\"uhrt nach einem Schritt den \"aquivalenten \"Ubergang.

\emph{Beweis per Induktion \"uber die Schrittanzahl:}

\emph{Basis ($t = 0$):} Anfangszustand $q_0^{\mathrm{ITM}} = q_0^T$ (gleich per Konstruktion).
Anfangsband: $\mathcal{B}_{\mathrm{ITM}}(i) = C_{\mathcal{B}_T(i)}$ per Kodierung. \checkmark

\emph{Induktionsschritt ($t \to t+1$):}
Sei zum Zeitpunkt $t$: $(q_t, \mathcal{B}_t^{\mathrm{ITM}}, k_t) \equiv (q_t, \mathcal{B}_t^T, k_t)$ (\"aquivalent).
TM-Schritt: $\delta_T(q_t, \mathcal{B}_t^T(k_t)) = (q_{t+1}, s', D)$.
ITM-Schritt: $\delta_{\mathrm{ITM}}(q_t, C_{\mathcal{B}_t^T(k_t)}) = (q_{t+1}, C_{s'}, D)$.
Neuer Bandzustand: $\mathcal{B}_{t+1}^{\mathrm{ITM}}(k_t) = C_{s'} = C_{\mathcal{B}_{t+1}^T(k_t)}$.
Neue Position: $k_{t+1}$ identisch (gleiche Bewegungsregel $D$).
Per Induktionshypothese sind alle anderen Bandzellen \"aquivalent.
Also $(q_{t+1}, \mathcal{B}_{t+1}^{\mathrm{ITM}}, k_{t+1}) \equiv (q_{t+1}, \mathcal{B}_{t+1}^T, k_{t+1})$. \checkmark

\textbf{Terminierung:}
$\mathcal{T}$ h\"alt in $q_{\mathrm{halt}}^T$ gdw.\ $\ITM$ in $q_{\mathrm{halt}}^{\mathrm{ITM}}$
h\"alt (per Konstruktion identische Haltebedingung).

\textbf{Ausgabekorrektheit:}
Die Ausgabe von $\ITM$ ist $\mathrm{dec}(\mathcal{B}_{\mathrm{final}})$:
Jede Zelle $C_s$ wird durch $s = C_s \cdot (|\Gamma_T|-1)/c_{\max}$ dekodiert
(inverse Kodierung). Per Induktionsschritt stimmt dies mit der TM-Ausgabe
\"uberein.

\textbf{Laufzeit:}
Jeder TM-Schritt wird in genau einem ITM-Schritt simuliert (Direktkodierung).
Damit $t_{\mathrm{ITM}}(n) = t_{\mathcal{T}}(n)$ (exakt, kein Overhead).

\textbf{Turing-\"Aquivalenz:}
$\mathcal{T} \leq_T \ITM$ (obige Simulation).
$\ITM \leq_T \mathcal{T}$ (Simulation der ITM durch eine TM ist trivial:
Die TM simuliert die ITM durch Kodierung der Konzentrationen als Bandzeichen;
jeder ITM-Schritt wird in $O(|\Gamma|)$ TM-Schritten simuliert).
Damit $\ITM \equiv_T \mathcal{T}$.
\end{proof}

\begin{korollar}[Universelle ITM]
\label{kor:uitm}
Es existiert eine \emph{universelle ITM} $\mathcal{T}_{\mathrm{ion},U}$, die bei Eingabe
$\langle \ITM, x \rangle$ die Berechnung von $\ITM$ auf $x$ simuliert.
\end{korollar}

\begin{proof}
Die universelle TM $\mathcal{T}_U$ existiert (klassisches Ergebnis~\cite{turing1936}).
Nach Satz~\ref{satz:itm_universal} gibt es eine ITM $\mathcal{T}_{\mathrm{ion},U}$, die $\mathcal{T}_U$
simuliert. Da $\mathcal{T}_U$ jede TM (und damit jede ITM) simuliert,
simuliert $\mathcal{T}_{\mathrm{ion},U}$ jede ITM.
\end{proof}

\begin{korollar}[Halteproblem der ITM]
\label{kor:halt}
Das \emph{ionotronische Halteproblem}
$\mathrm{HALT}_{\mathrm{ITM}} = \{(\langle \ITM \rangle, x) \mid \ITM \text{ h\"alt auf } x\}$
ist unentscheidbar.
\end{korollar}

\begin{proof}
Angenommen, es gibt einen Entscheider $D$ f\"ur $\mathrm{HALT}_{\mathrm{ITM}}$.
Dann gibt es (durch Simulation via Satz~\ref{satz:itm_universal}) auch
einen Entscheider f\"ur $\mathrm{HALT}_{\mathrm{TM}}$.
Dies widerspricht dem klassischen Unentscheidbarkeits-Satz von Turing~\cite{turing1936}.
\end{proof}

\begin{figure}[h!]
\centering
\includegraphics[width=0.95\textwidth]{plot6_universalitaet.pdf}
\caption{Berechnungsuniversalit\"at der ITM.
Oben links: Busy-Beaver $BB(2)$ Simulation auf der ITM.
Oben rechts: Berechnungsklassen-Hierarchie.
Unten: NF-2/NF-6-Verbindungen und Pareto-Analyse.}
\label{fig:universalitaet}
\end{figure}

%%%%%%%%%%%%%%%%%%%%%%%%%%%%%%%%%%%
\section{Zusammenfassung}
%%%%%%%%%%%%%%%%%%%%%%%%%%%%%%%%%%%

Diese Arbeit hat die vollst\"andige formale Theorie der Ionotronischen
Turing-Maschine entwickelt. Die wichtigsten Ergebnisse:

\begin{itemize}
    \item \textbf{Lemma~\ref{lem:kappa_mono}} (Kohlrausch-Monotonie):
          $\kappa(c)$ ist f\"ur $c \leq c^*$ streng monoton; Beweis via Differentialrechnung.
    \item \textbf{Lemma~\ref{lem:donnan}} (Donnan-Gleichgewicht):
          $E_D = (RT/zF)\ln(a_i^\beta/a_i^\alpha)$; vollst\"andiger thermodynamischer Beweis.
    \item \textbf{Lemma~\ref{lem:inject}} (Ionische Kodierungs-Injektivit\"at):
          Rauschstabilit\"at mit Fehlerrate $\leq 2e^{-\Delta^2/(2\sigma^2)}$.
    \item \textbf{Satz~\ref{satz:kohlrausch}} (Kohlrausch-Schaltungssatz):
          AND, OR, NOT ionisch realisiert; vollst\"andige Fallanalyse aller 4 Eingangsf\"alle.
    \item \textbf{Satz~\ref{satz:vollstaendig}} (Funktionale Vollst\"andigkeit):
          $\{\mathrm{AND}, \mathrm{NOT}\}$ vollst\"andig; Beweis via DNF.
    \item \textbf{Korollar~\ref{kor:nand}} (NAND-Vollst\"andigkeit).
    \item \textbf{Satz~\ref{satz:phys_real}} (Physikalische Realisierbarkeit):
          Explizite Berechnungen aller Zeitkonstanten.
    \item \textbf{Satz~\ref{satz:analog_info}} (Analog-Information):
          $I = \log_2(\mathrm{DR}) \cdot \ell/\ell_{\mathrm{cell}}$ bit.
    \item \textbf{Satz~\ref{satz:diskr_verlust}} (Diskretisierungsverlust):
          SQNR-Formel $= 6{,}02n + 1{,}76\,\mathrm{dB}$.
    \item \textbf{Lemma~\ref{lem:linear_shift}} (LSPR-Linearit\"at):
          $\Delta\lambda = S \cdot \Delta n + O(\Delta n^2)$.
    \item \textbf{Satz~\ref{satz:hybrid}} (Hybrid-Satz):
          $\Delta\lambda = 0{,}24\,c\,\mathrm{nm}$ mit Fehler $< 1\%$.
    \item \textbf{Satz~\ref{satz:itm_universal}} (Turing-Universalit\"at):
          Vollst\"andiger Induktionsbeweis der Turing-\"Aquivalenz ITM $\equiv_T$ TM.
    \item \textbf{Korollar~\ref{kor:halt}} (Unentscheidbarkeit):
          Halteproblem der ITM ist unentscheidbar.
\end{itemize}

%%%%%%%%%%%%%%%%%%%%%%%%%%%%%%%%%%%
\section{Ausblick}
%%%%%%%%%%%%%%%%%%%%%%%%%%%%%%%%%%%

\subsection{Biologische Neuronen als ITM: Hodgkin-Huxley-Formalisierung}

Das Hodgkin-Huxley-Modell~\cite{hodgkin1952} beschreibt das Aktionspotential
eines Neurons durch ein System gew\"ohnlicher Differentialgleichungen:
\begin{align*}
C_m \frac{dV}{dt} &= -g_{\mathrm{Na}} m^3 h (V - E_{\mathrm{Na}})
                     -g_K n^4 (V - E_K) - g_L (V - E_L) + I_{\mathrm{ext}},\\
\frac{dm}{dt} &= \alpha_m(V)(1-m) - \beta_m(V)m, \quad \text{(analog f\"ur } n, h\text{)},
\end{align*}
wobei $m, n, h \in [0,1]$ Aktivierungs- und Inaktivierungsgating-Variablen sind.
Diese Gleichungen beschreiben exakt das ionotronische Rechenprinzip:
Na$^+$- und K$^+$-Kan\"ale kodieren, verarbeiten und \"ubertragen Information.
Formale Verbindung zur ITM: Der \"Ubergang $\delta(q, b) = (q', b', D)$
entspricht dem Spike-Muster des Neurons;
Bandbewegung entspricht der Signalweiterleitung entlang des Axons.
Offene Frage: Ist ein biologisches neuronales Netz Turing-vollst\"andig?
McCulloch-Pitts-Neuronen~\cite{mcculloch1943} sind es -- die ITM
liefert die physikalische Fundierung.

\subsection{Ionotronische FPGA-Architektur}

Eine \emph{ionotronische FPGA-Architektur} nutzt rekonfigurierbare
Mikrokanal-Logikblöcke:
Jeder Block enth\"alt eine $2 \times 2$ Mikrokanalmatrix, die durch Konzentrations\"anderung
zwischen AND, OR, NOT, XOR konfigurierbar ist.
Verbindungen werden durch Gate-Kan\"ale realisiert (\"ahnlich CMOS-Pass-Transistoren).
Vorteil gegen\"uber SRAM-FPGAs: nicht-fl\"uchtige Konfigurationsspeicherung
(ionische Kapazit\"at h\"alt Konfiguration auch ohne Versorgungsspannung),
biokompatibel, l\"oslich.
Verbindung zu~\cite{epp2026fpga}: Der Subgraph Algorithmus~\cite{epp2026subgraph}
optimiert die ionotronische FPGA-Topologie in $O(n^3)$.

\subsection{Quantenionotronik}

F\"ur Ionen in Ionenfallen bei $T < 1\,\mu$K~\cite{cirac1995} sind
Quantenkoh\"arenz und ionisches Rechnen kombinierbar.
Quantensuperposition von Konzentrationslevels:
$|\psi\rangle = \alpha|c_{\mathrm{low}}\rangle + \beta|c_{\mathrm{high}}\rangle$
(koherte Superposition von Bit~0 und Bit~1 in einem Quanten-Ionenkanal).
Diese Quanten-ITM (QITM) erg\"anzt die ITM um Quanten-Parallelismus.
Verbindung zu~\cite{epp2026qecc}: Surface-Code-Fehlerkorrektur f\"ur QITM.

\subsection{Neuromorphes Rechnen und ionotronische Synapsen}

Biologische Synapsen sind ionotronische Schnittstellen.
Ein ionotronisches Synapse-Modell:
Synapsenstärke $w_{ij}$ $\leftrightarrow$ Kanal-Leitfähigkeit $\kappa_{ij}$.
Hebbian Learning: $\Delta w_{ij} = \eta\,x_i\,x_j$
$\leftrightarrow$ $\Delta\kappa_{ij} = \eta\,c_i\,c_j$.
Spike-Timing-Dependent Plasticity (STDP):
$\Delta w = A_+\,e^{-\Delta t/\tau_+}$ f\"ur pr\"asynapsoide Aktivierung
entspricht einer ionischen Diffusionskinetik.
Verbindung zu~\cite{epp2026nf6}: Kochlea-Architektur als frequenzselektives
ionotronisches neuronales Netz.

\subsection{Ionotronische Kryptographie: Diffusionsbasierte Verschl\"usselung}

Fick'sche Diffusion ist physikalisch irreversibel (Entropieproduktion).
Ionotronische Verschl\"usselung: Schlüssel $K$ ist das Konzentrationsprofil
$\{c_i(0)\}_{i=1}^n$ zum Zeitpunkt $t=0$.
Ciphertext: $\{c_i(\Delta t)\}_{i=1}^n$ nach Diffusionszeit $\Delta t$
gem\"a{\ss} $\partial c/\partial t = D\nabla^2 c$.
Entschl\"usselung: L\"osung der inversen W\"armegleichung (ill-posed, instabil
ohne Schl\"ussel).
Verbindung zur Signatur-Chiffre~\cite{epp2026sigchiffre}:
Graphen-Signatur des Konzentrationsprofils als kryptographischer Schl\"ussel.

\subsection{Lab-on-Chip: In-situ-Diagnose mit ITM}

Ein Lab-on-Chip integriert biologische Probenaufbereitung und
Informationsverarbeitung auf einem Chip-Format.
Ionotronische Integration: Das Patientenblut oder Wasser der Probe
dient direkt als Rechenmedium; Ionenkonzentrationen (Na$^+$, K$^+$, Ca$^2+$)
werden gemessen und durch die ITM-Band-Kodierung sofort als Information
verfügbar.
Verbindung zu~\cite{epp2026norovirus}: Norovirus-Epitop-Bindung \"andert
lokale $\zeta$-Potential und damit $\kappa$; Subgraph Algorithmus
erkennt Virussignatur in $O(n^3)$.

\subsection{Ionotronische Simulation partieller Differentialgleichungen}

Ionotronische Systeme sind nat\"urliche Analogrechner f\"ur PDEs:
Die Fick'sche Diffusionsgleichung $\partial c/\partial t = D\nabla^2 c$
l\"ost die W\"armeleitungsgleichung \emph{exakt} (physikalisch, ohne Diskretisierungsfehler).
Erweitert: Nernst-Planck-Gleichung
$\partial c/\partial t = D\nabla^2 c - \mu_{\mathrm{ion}} \nabla(\mathbf{E}c)$
modelliert drift-diffusion-Gleichungen (Halbleiterphysik, Neuronales Modellieren,
atmosph\"arische Transportprozesse).
Verbindung zum Analog-Paradigma~\cite{epp2026analog}: Ionotronische
PDE-Simulation ist die physikalische Instantiierung des Analog-Rechners.

\subsection{Skalierung: Von Nano zu Mikro zu Makro}

\paragraph{Nano-ITM ($d < 10\,\text{nm}$).}
Im Angstrom-Regime zeigen Ionenkan\"ale Quanteneffekte~\cite{bocquet2020}:
Einzel-Ionen-Tunneln, diskrete Leitf\"ahigkeitssprünge.
Eine Nano-ITM mit Kohlenstoffnanor\"ohren ($d \approx 1\,\text{nm}$) h\"atte
$V_{\mathrm{cell}} \approx 10^{-24}\,$L und $\sigma_c \approx 0{,}01\,$mol/L
(extrem klein, aber diskrete Quantenfluktuationen dominant).

\paragraph{Mikro-ITM ($d = 1$--$100\,\mu$m).}
Das in dieser Arbeit behandelte Regime.
Optimales Betriebsfenster: $c \in [10^{-3}, 1]\,\text{mol/L}$,
$V_{\mathrm{cell}} \in [10^{-15}, 10^{-9}]\,\text{L}$.

\paragraph{Makro-ITM ($d > 1\,\text{mm}$).}
Flusssysteme, Mikroreaktoren, Bio-Fermenter als ``Computer''.
Anwendung: Stoffwechselsteuerung in biomedizinischen Anlagen.

\subsection{Verbindung zur Signalverarbeitung und Fourier-Analyse}

Kochlea-Geometrie als physikalisches Spektralanalyseger\"at:
Die abnehmende Kanalbreite entlang der Kochlea realisiert einen
\"ortlich aufgel\"osten Bandpassfilter -- ein physikalisches Fourier-Spektrum.
Formal: Das Konzentrationsprofil $c(x)$ hat eine Fourier-Darstellung
$\hat{c}(k) = \int_0^L c(x) e^{-ikx} dx$, die durch ionische Diffusion
nat\"urlich in niedrig-$k$- und hoch-$k$-Komponenten zerlegt wird.
Verbindung zur Signaltheorie~\cite{epp2026signalth}:
Ionotronische Verarbeitung ist ein \emph{analoges Fourier-Spektrometer},
das ohne Diskretisierung (und damit verlustfrei) arbeitet.

\subsection{Umweltvertr\"aglichkeit und Nachhaltigkeit}

\paragraph{Rohstoffe.}
Wasser und NaCl: abundante, ungiftige, biologisch abbaubare Materialien.
Im Gegensatz zu CMOS-Herstellung (Flusss\"aure-\"Atzung, Borphosphor-Silikatglas,
seltene Erden) ist die ITM-Fabrikation ohne Gefahrstoffe m\"oglich.

\paragraph{Energie.}
Ionotronisches Rechnen bei physiologischen Konzentrationen:
$E_{\mathrm{ion}} \approx k_BT \approx 4 \cdot 10^{-21}\,$J/Operation.
CMOS-Minimum (ITRS 2025): $E_{\mathrm{CMOS}} \approx 10^{-18}\,\text{J/op}$.
Ratio: $E_{\mathrm{ion}}/E_{\mathrm{CMOS}} \approx 4 \cdot 10^{-3}$
-- Ionotronik ist \emph{drei Gr\"o{\ss}enordnungen} energieeffizienter.

\paragraph{Entsorgung.}
Ein ionotronischer Chip l\"ost sich in Wasser auf (biodegradabel).
Keine elektronische Sonderm\"ullentsorgung notwendig.

\paragraph{Meeresumgebungen.}
Verbindung zu~\cite{epp2026cleanocn}: Meerwasser als Rechenmedium
f\"ur Unterwasser-Datenverarbeitungseinheiten, ohne Frischwasserversorgung.

%%%%%%%%%%%%%%%%%%%%%%%%%%%%%%%%%%%
\begin{thebibliography}{99}
%%%%%%%%%%%%%%%%%%%%%%%%%%%%%%%%%%%

\bibitem{turing1936}
Turing, A.\,M. (1936).
On computable numbers, with an application to the Entscheidungsproblem.
\textit{Proc.\ London Math.\ Soc.}, 42(2), 230--265.

\bibitem{shockley1948}
Shockley, W., Bardeen, J., \& Brattain, W.\,H. (1948).
The transistor, a semi-conductor triode.
\textit{Physical Review}, 74(2), 232--233.

\bibitem{dennard1974}
Dennard, R.\,H. et al.\ (1974).
Design of ion-implanted MOSFET's with very small physical dimensions.
\textit{IEEE J.\ Solid-State Circuits}, 9(5), 256--268.

\bibitem{haensch2006}
Haensch, W. et al.\ (2006).
Silicon CMOS devices beyond scaling.
\textit{IBM Journal of Research and Development}, 50(4), 339--361.

\bibitem{dongarra2011}
Dongarra, J.\ et al.\ (2011).
The International Exascale Software Project roadmap.
\textit{IJHPCA}, 25(1), 3--60.

\bibitem{yang2018}
Yang, C.\ et al.\ (2018).
Ionic rectification and memristive behavior in soft ionotronic channels.
\textit{Advanced Materials}, 30(14), 1705708.

\bibitem{robin2021}
Robin, P.\ et al.\ (2021).
Long-term memory and synapse-like dynamics in two-dimensional nanofluidic channels.
\textit{Science}, 373(6555), 687--691.

\bibitem{bocquet2020}
Bocquet, L.\ (2020).
Nanofluidics coming of age.
\textit{Nature Materials}, 19, 254--256.

\bibitem{hodgkin1952}
Hodgkin, A.\,L., \& Huxley, A.\,F. (1952).
A quantitative description of membrane current.
\textit{J.\ Physiology}, 117(4), 500--544.

\bibitem{mcculloch1943}
McCulloch, W.\,S., \& Pitts, W.\ (1943).
A logical calculus of ideas immanent in nervous activity.
\textit{Bull.\ Math.\ Biophysics}, 5(4), 115--133.

\bibitem{markram2006}
Markram, H.\ (2006).
The Blue Brain Project.
\textit{Nature Reviews Neuroscience}, 7(2), 153--160.

\bibitem{shannon1948}
Shannon, C.\,E.\ (1948).
A mathematical theory of communication.
\textit{Bell System Technical Journal}, 27, 379--423.

\bibitem{sipser2013}
Sipser, M.\ (2013).
\textit{Introduction to the Theory of Computation}, 3rd ed.
Cengage Learning.

\bibitem{cirac1995}
Cirac, J.\,I., \& Zoller, P.\ (1995).
Quantum computations with cold trapped ions.
\textit{Physical Review Letters}, 74(20), 4091--4094.

\bibitem{maier2007}
Maier, S.\,A.\ (2007).
\textit{Plasmonics: Fundamentals and Applications}.
Springer.

\bibitem{laser2004}
Laser, D.\,J., \& Santiago, J.\,G.\ (2004).
A review of micropumps.
\textit{J.\ Micromechanics and Microengineering}, 14(6), R35--R64.

\bibitem{lisdat2008}
Lisdat, F., \& Sch\"afer, D.\ (2008).
The use of electrochemical impedance spectroscopy for biosensing.
\textit{Anal.\ Bioanal.\ Chem.}, 391(5), 1555--1567.

\bibitem{quan2009}
Quan, X., Fang, N., \& Scriven, L.\,E.\ (2009).
Refractive index of NaCl solutions.
\textit{J.\ Chemical \& Engineering Data}, 54(2), 642--645.

\bibitem{epp2026iono}
Epp, S.\ (2026).
Ionotronisches Fl\"ussigkeits-Rechnen (IRF).
Universit\"at Bielefeld.
\url{https://github.com/hjstephan86/science/tree/main/science/iono}

\bibitem{epp2026oqtum}
Epp, S.\ (2026).
Aktive Nanooptik -- programmierbare photonische Oberfl\"achen (oqtum).
Universit\"at Bielefeld.
\url{https://github.com/hjstephan86/science/tree/main/science/oqtum}

\bibitem{epp2026ionoxoqtum}
Epp, S.\ (2026).
Effiziente Nanostrukturen mit Wasser und Licht (ionoxoqtum).
Universit\"at Bielefeld.
\url{https://github.com/hjstephan86/science/tree/main/science/ionoxoqtum}

\bibitem{epp2026analog}
Epp, S.\ (2026).
Analog als prim\"arer Begriff zur Beschreibung der realen Welt.
Google Drive (epp-group).

\bibitem{epp2026digi}
Epp, S.\ (2026).
Von der Diskretion.
Universit\"at Bielefeld.
\url{https://github.com/hjstephan86/science/tree/main/science/digi}

\bibitem{epp2026nf_all}
Epp, S.\ (2026).
Neue Forschungsfelder (nf\_all) -- NF-1 bis NF-8.
Google Drive (epp-group).

\bibitem{epp2026subgraph}
Epp, S.\ (2026).
Der Subgraph Algorithmus -- $O(n^3)$.
Universit\"at Bielefeld.
\url{https://github.com/hjstephan86/science}

\bibitem{epp2026fpga}
Epp, S.\ (2026).
Subgraph-basierte Topologie-Optimierung von FPGA-Beschleunigern.
Google Drive (epp-group).

\bibitem{epp2026qecc}
Epp, S.\ (2026).
Topologische Quantenfehlerkorrektur.
Google Drive (epp-group).

\bibitem{epp2026sigchiffre}
Epp, S.\ (2026).
Signatur-Chiffre.
Universit\"at Bielefeld.
\url{https://github.com/hjstephan86/science/tree/main/science/sigchiffre}

\bibitem{epp2026norovirus}
Epp, S.\ (2026).
Norovirus-Impfstoff-Entwicklung via Subgraph-Analyse.
Google Drive (epp-group).

\bibitem{epp2026nf6}
Epp, S.\ (2026).
NF-6: Biomimetische Kochlea-Architektur.
Google Drive (epp-group).

\bibitem{epp2026signalth}
Epp, S.\ (2026).
Signaltheorie -- Das Wunder der e-Funktion.
Universit\"at Bielefeld.
\url{https://github.com/hjstephan86/science/tree/main/science/signalth}

\bibitem{epp2026cleanocn}
Epp, S.\ (2026).
Systematische Analyse von Ozean-Reinigungsstrategien.
Universit\"at Bielefeld.
\url{https://github.com/hjstephan86/science/tree/main/science/cleanocn}

\end{thebibliography}

\end{document}
